\chapter{Abstract}

\begin{tabular}{@{} >{\bfseries}p{0.22\textwidth} p{0.75\textwidth} @{}}
\MakeUppercase{Keywords} &
Inverse Estimation, Adaptive Overcurrent Protection, Synchronous Generator,
Transformer Differential, IBR Line Differential, Levenberg--Marquardt,
Confidence Scoring, Digital Twin, Physics-Informed Neural Network,
Wide-Area Protection Coordination, IEC~61850, Adversarial Robustness.
\end{tabular}
\par

\begin{doublespacing}

The global power grid is undergoing a fundamental structural transformation:
synchronous machines---whose high-inertia fault signatures underpinned a century
of deterministic relay design---are being displaced at accelerating pace by
Inverter-Based Resources (IBRs). Grid-Following (GFL) and Grid-Forming (GFM)
converters are software-controlled current sources whose fault currents are
limited to 1.1--2.0~pu by fast sub-cycle current limiters, and whose sequence
contributions are determined by firmware rather than physical machine reactance.
This paradigm shift systematically invalidates the foundational assumptions of
legacy overcurrent, distance, and differential relay algorithms.

This thesis proposes the \textbf{Self-Adaptive Model-Based Protection (SAMBP)}
architecture for future power systems. The framework treats adaptive protection
as an inverse estimation problem: each protection element maintains an explicit
physical model of its protected device and identifies model parameters from
real-time measurements by the Levenberg--Marquardt algorithm. The framework
comprises three tiers and 25 original contributions (17 validated, 8 pending
an 18--24 month experimental execution phase).

Chapter~4 develops five adaptive 87L line differential algorithms achieving
TPR~= 1.000 across $k_\mathrm{ibr} \in [0.03, 0.85]$~pu, including 87LN and
87LN0 elements for pure-IBR SLG coverage. Chapter~5 introduces SAMBP-SyncOC:
the first inverse-estimation-based adaptive overcurrent protection for
synchronous generators, achieving $R^2 > 0.999$ and 60\% reduction in
spurious relay pickups. Chapter~6 develops the 87T-SPRT transformer
differential scheme using Wald sequential probability ratio test inrush
discrimination with explicit false-trip and miss-rate guarantees.
Chapters~7--8 validate quadrilateral distance protection and GFL/GFM
mode-switching islanding. Chapter~9 delivers a non-blocking EKF digital twin
(RMSE $\leq 0.024$~pu). Chapter~10 presents a physics-informed fault
classifier (96.1\% accuracy, zero physics violations). Chapter~11 integrates
all components into the Wide-Area Protection Coordinator. Chapter~12 establishes
the cybersecurity framework including adversarial robustness against CT-injection
attacks.

The SAMBP framework is the first published architecture to integrate model-based
relay adaptation, digital-twin state estimation, physics-informed classification,
and closed-loop coordination across all protection elements with provable safety
guarantees.

\end{doublespacing}
